\documentclass[11pt]{article}

\usepackage[margin=1in]{geometry}
\usepackage{amsmath,amssymb}
\usepackage{booktabs}
\usepackage{longtable}
\usepackage{array}
\usepackage{enumitem}
\usepackage{xcolor}
\usepackage{hyperref}

\hypersetup{
  colorlinks=true,
  linkcolor=blue,
  citecolor=blue,
  urlcolor=blue
}

\newcommand{\Ne}{N_{\mathrm e}}
\newcommand{\Te}{T_{\mathrm e}}
\newcommand{\Ti}{T_{\mathrm i}}
\newcommand{\Vi}{V_{\mathrm i}}
\newcommand{\vlos}{v_{\mathrm{los}}}
\newcommand{\R}{\mathbb{R}}
\newcommand{\E}{\mathbb{E}}
\newcommand{\dd}{\mathrm{d}}
\newcommand{\kmps}{\mathrm{km\,s^{-1}}}
\newcommand{\mps}{\mathrm{m\,s^{-1}}}

\title{Research Plan for Uncertainty-Aware INR Reconstruction of Multi-Beam ISR Plasma Fields}
\author{Prepared for the UAI PostDoc INR--Radar Project}
\date{July 2026}

\begin{document}

\maketitle

\begin{abstract}
This document consolidates the current state of the implicit neural representation (INR) radar
project, the synthetic-data results obtained so far, the proposal commitments, and a recommended
research plan for the next stage. The central objective is to reconstruct continuous ionospheric
plasma fields from sparse multi-beam incoherent scatter radar (ISR) products while providing
explicit reliability maps. The proposed path remains deliberately outside a full physics-informed
neural network (PINN): physical knowledge enters as uncertainty-aware likelihoods, anisotropic
soft priors, geometry-aware support diagnostics, integration-time-aware synthetic tests, and
calibrated reliability products. A key addition in this version is that integration time must be
treated as a first-order experimental axis. For fast plasma motion, one-minute, two-minute,
five-minute, and ten-minute products are not merely different cadences; they represent different
time-averaged measurements of a moving field and can therefore contain motion blur and
stationarity errors.
\end{abstract}

\tableofcontents

\section{Executive Summary}

The project should be framed as an uncertainty-aware continuous interpolation and reconstruction
tool for ISR-derived plasma parameters. The tool should produce two coupled outputs:
\begin{enumerate}[leftmargin=*]
  \item a reconstructed field, initially $\log_{10}\Ne(x,y,t)$ and later
  $\log_{10}\Ne(x,y,z,t)$, $\Te(x,y,z,t)$, and $\Ti(x,y,z,t)$;
  \item a reliability layer that tells users where the reconstruction is data constrained,
  weakly constrained, or prior dominated.
\end{enumerate}

The work so far supports the following conclusions.
\begin{enumerate}[leftmargin=*]
  \item A SIREN-like INR is a good baseline because it provides a differentiable continuous field
  and meaningful derivatives \cite{Sitzmann2020SIREN}.
  \item A pure data-only INR can fit sparse observations but may produce unsupported structure
  and derivative artifacts between beams.
  \item Moderate curvature regularization improved the synthetic reconstruction in the first
  controlled sweep, while stronger regularization was not uniformly beneficial.
  \item The current four-case regularization sweep is a pilot, not final evidence. A defensible
  study requires a structured loss-weight protocol, multiple velocities, multiple integration
  times, multiple beam counts, and withheld-beam/time validation.
  \item Integration time must be included. At $2~\kmps$, a plasma feature moves $120~\mathrm{km}$
  during a one-minute integration, $240~\mathrm{km}$ during two minutes, $600~\mathrm{km}$
  during five minutes, and $1200~\mathrm{km}$ during ten minutes. These are not small
  perturbations for the current synthetic domains.
  \item Eleven-beam experiments should not be discarded. They are likely too sparse for uniformly
  accurate reconstruction, but they are valuable stress tests for reliability maps. A successful
  method should warn the user when the eleven-beam reconstruction is not trustworthy.
\end{enumerate}

The recommended next milestone is therefore an integration-aware synthetic benchmark:
\begin{enumerate}[leftmargin=*]
  \item beam counts: 42, 23, and 11;
  \item velocities: $0$, $0.36$, $0.75$, $1.0$, $1.5$, $2.0$, and $3.0~\kmps$;
  \item integration times: 1, 2, 5, and 10 minutes;
  \item model families: baseline INR in $(x,y,t)$, velocity-aware co-moving INR, and later
  full 4D INR;
  \item reliability methods: MC dropout plus geometry/support features, and conformal prediction;
  \item validation: dense synthetic truth, withheld beams, withheld times, regional metrics, and
  interval coverage.
\end{enumerate}

\section{Project Contract from the Proposal}

The UAI proposal frames the project as an uncertainty-aware reconstruction method for 4D
ionospheric structure from sparse multi-beam radar data \cite{Diaz2026Proposal}. The core
commitments are:
\begin{enumerate}[leftmargin=*]
  \item fit continuous neural fields for ISR-derived parameters such as $\Ne$, $\Te$, and $\Ti$;
  \item use ISR-reported errors as heteroscedastic measurement noise;
  \item add two interpretable soft priors: temporal continuity and anisotropic smoothness;
  \item deliver reliability products that include an observability proxy;
  \item validate using withheld-beam and withheld-time tests;
  \item benchmark against interpolation and smoothing baselines.
\end{enumerate}

One update is recommended. The proposal mentions ensemble-style uncertainty in some places,
but the current plan should not use ensembles as a main method. The uncertainty branch should be:
\begin{enumerate}[leftmargin=*]
  \item MC dropout plus geometry/support features \cite{Gal2016Dropout,KendallGal2017};
  \item conformal prediction, especially normalized or group-conditioned conformal intervals
  \cite{AngelopoulosBates2021,Romano2019CQR}.
\end{enumerate}

The revised project claim should be:
\begin{quote}
We reconstruct continuous ISR-derived plasma fields and provide calibrated reliability diagnostics
that distinguish data-constrained interpolation from prior-dominated extrapolation.
\end{quote}

\section{Current Methodological Progress}

\subsection{Data Representation}

The current real-data branch uses sparse radar samples as coordinate-value pairs,
\begin{equation}
  \mathcal{D} = \{(x_i,y_i,t_i,u_i)\}_{i=1}^{N},
  \qquad
  u_i = \log_{10}\Ne{}_{i},
\end{equation}
where $(x,y)$ are horizontal coordinates, $t$ is time, and $u$ is log electron density.
Coordinates and targets are normalized during training. The model is a coordinate-based neural
field,
\begin{equation}
  f_{\theta}(x,y,t) \approx \log_{10}\Ne(x,y,t).
\end{equation}

The current implementation is not yet full $(x,y,z,t)$. It uses altitude-filtered slices, which is
appropriate at this stage because it keeps the validation problem interpretable.

\subsection{Baseline Architecture}

The current baseline is a sine-activated MLP inspired by SIREN:
\begin{equation}
  h_{\ell} = \sin\left(\omega_{\ell}(W_{\ell}h_{\ell-1}+b_{\ell})\right).
\end{equation}
SIREN is attractive because it represents continuous fields and derivatives well
\cite{Sitzmann2020SIREN}. This matters because the project evaluates first and second derivatives
for regularization and diagnostics. A ReLU network would be poorly matched to second-derivative
curvature penalties because its second derivatives are zero almost everywhere and undefined at
kinks.

\subsection{Current Loss}

The current training objective is:
\begin{equation}
  \mathcal{L}
  =
  \mathcal{L}_{\mathrm{data}}
  +
  \lambda_{xy}\mathcal{L}_{xy}
  +
  \lambda_t\mathcal{L}_t.
\end{equation}
The data loss is currently MSE or synthetic MSE:
\begin{equation}
  \mathcal{L}_{\mathrm{data}}
  =
  \frac{1}{N}\sum_i
  \left(f_{\theta}(x_i,y_i,t_i)-u_i\right)^2.
\end{equation}

The spatial term is a horizontal Hessian penalty:
\begin{equation}
  \mathcal{L}_{xy}
  =
  \E_{(x,y,t)\sim\Omega}
  \left[
    f_{xx}^2 + 2f_{xy}^2 + f_{yy}^2
  \right].
\end{equation}
This is not the same as a Laplacian penalty $(f_{xx}+f_{yy})^2$. It penalizes curvature in all
horizontal directions, including mixed curvature.

The temporal term is:
\begin{equation}
  \mathcal{L}_{t}
  =
  \E_{(x,y,t)\sim\Omega}
  \left[
    f_{tt}^2
  \right].
\end{equation}
This allows linear temporal trends but discourages rapid temporal bending.

\subsection{Synthetic Validation Infrastructure}

The synthetic branch was created because real radar data do not provide dense truth. The synthetic
case allows the same INR training workflow to be tested against a known analytic plasma field.
The completed progress includes:
\begin{enumerate}[leftmargin=*]
  \item synthetic moving-patch generator;
  \item sparse radar-like observation geometries;
  \item 42-beam and 23-beam sampling patterns;
  \item data-only and regularized training cases;
  \item derivative diagnostics even when regularization is disabled;
  \item dense reconstruction analysis using the corrected reconstruction/error script;
  \item regional error metrics: full domain, interior, near-observation, and high-gradient regions;
  \item comparison plots and summary tables.
\end{enumerate}

\subsection{First Synthetic Sweep Results}

The first controlled sweep used a moving Gaussian-like density patch at approximately
$360~\mps$. The key cases were:
\begin{enumerate}[leftmargin=*]
  \item data-only;
  \item moderate regularization: $\rho_{xy}=0.30$, $\rho_t=0.30$;
  \item stronger spatial regularization: $\rho_{xy}=0.70$, $\rho_t=0.30$;
  \item stronger spatial and temporal regularization: $\rho_{xy}=0.70$, $\rho_t=0.70$.
\end{enumerate}

The main finding was that moderate regularization was best overall in the first benchmark. It
reduced artifacts and improved gradient recovery without over-smoothing as strongly as the higher
regularization cases. The data-only run was competitive for scalar density in some high-gradient
regions, but derivative recovery and artifact control were weaker. This supports the project
hypothesis that priors are useful, but it does not yet prove that any particular weight is optimal.

\section{Why Integration Time Must Be Included}

\subsection{Cadence and Integration Time Are Different Concepts}

ISR products commonly appear at integration times such as 1, 2, 5, and 10 minutes. It is tempting
to treat these products as point samples at their timestamps. That is only approximately true when
the plasma is slowly varying over the integration interval. For fast plasma flow, a radar product is
better interpreted as a time-averaged measurement over a moving plasma field.

Two time scales must be separated:
\begin{enumerate}[leftmargin=*]
  \item sampling interval or cadence, $\Delta t_{\mathrm{samp}}$: the time between successive records;
  \item integration duration, $T_{\mathrm{int}}$: the time over which the measurement is accumulated.
\end{enumerate}

Often these are similar in processed products, but conceptually they have different effects.
Cadence controls the displacement between frames. Integration duration controls motion blur within
each frame.

\subsection{Forward Model for Integrated Observations}

An instantaneous synthetic observation would be:
\begin{equation}
  y_i = u(\mathbf{r}_i,t_i) + \epsilon_i,
\end{equation}
where $\mathbf{r}_i=(x_i,y_i,z_i)$ or $(x_i,y_i)$.

An integration-aware observation should instead be:
\begin{equation}
  y_i
  =
  \frac{1}{T_{\mathrm{int}}}
  \int_{t_i-T_{\mathrm{int}}/2}^{t_i+T_{\mathrm{int}}/2}
  u(\mathbf{r}_i,\tau)\,\dd \tau
  +
  \epsilon_i.
  \label{eq:integrated_observation}
\end{equation}

If the plasma field is a moving structure,
\begin{equation}
  u(\mathbf{r},t) = u_0(\mathbf{r}-\mathbf{v}t),
\end{equation}
then Eq.~\ref{eq:integrated_observation} is a convolution of the instantaneous structure with a
line segment in the direction of motion. The measurement is not just noisy; it is blurred along
the flow direction.

\subsection{Motion Blur for a Gaussian Patch}

For a Gaussian patch with horizontal covariance $\Sigma_0$ moving at velocity $\mathbf{v}$, a
uniform integration interval approximately adds covariance along the direction of motion:
\begin{equation}
  \Sigma_{\mathrm{eff}}
  \approx
  \Sigma_0
  +
  \frac{T_{\mathrm{int}}^2}{12}
  \mathbf{v}\mathbf{v}^{\top}.
  \label{eq:motion_blur_covariance}
\end{equation}
This is the covariance of a Gaussian convolved with a uniform line segment. The important
consequence is that motion blur is anisotropic and grows quadratically with integration time.

\subsection{Dimensionless Severity Numbers}

The synthetic benchmark should track dimensionless numbers rather than only raw velocities.
Let $\ell$ be a characteristic spatial scale, such as patch width, beam spacing, or the smallest
feature scale we expect to reconstruct. Define:
\begin{align}
  C_{\mathrm{samp}} &= \frac{\|\mathbf{v}\|\Delta t_{\mathrm{samp}}}{\ell},
  \\
  C_{\mathrm{int}} &= \frac{\|\mathbf{v}\|T_{\mathrm{int}}}{\ell}.
\end{align}

Interpretation:
\begin{enumerate}[leftmargin=*]
  \item $C_{\mathrm{samp}}\ll 1$: the feature moves little between records;
  \item $C_{\mathrm{samp}}\sim 1$: the feature moves by about one feature width between records;
  \item $C_{\mathrm{samp}}>1$: temporal interpolation becomes a tracking problem;
  \item $C_{\mathrm{int}}\ll 1$: the integrated product is close to instantaneous;
  \item $C_{\mathrm{int}}\sim 1$: the measurement is visibly motion blurred;
  \item $C_{\mathrm{int}}>1$: a point-sample interpretation is physically questionable.
\end{enumerate}

\subsection{Displacements for Candidate Velocities and Integration Times}

Table~\ref{tab:velocity_displacement} shows how far a structure moves during common ISR
integration intervals. These values strongly argue for including integration time in the synthetic
benchmark.

\begin{longtable}{rrrrr}
\caption{Horizontal displacement during integration for candidate synthetic velocities.}
\label{tab:velocity_displacement}\\
\toprule
Velocity & 1 min & 2 min & 5 min & 10 min \\
$(\kmps)$ & km & km & km & km \\
\midrule
\endfirsthead
\toprule
Velocity & 1 min & 2 min & 5 min & 10 min \\
$(\kmps)$ & km & km & km & km \\
\midrule
\endhead
0.00 & 0 & 0 & 0 & 0 \\
0.36 & 21.6 & 43.2 & 108 & 216 \\
0.75 & 45 & 90 & 225 & 450 \\
1.00 & 60 & 120 & 300 & 600 \\
1.50 & 90 & 180 & 450 & 900 \\
2.00 & 120 & 240 & 600 & 1200 \\
3.00 & 180 & 360 & 900 & 1800 \\
\bottomrule
\end{longtable}

At $2~\kmps$, a 5-minute product averages over roughly $600~\mathrm{km}$ of motion. For many
radar volumes this is not a detail; it changes the inverse problem. A 10-minute product at
active-time velocities may be excellent for reducing stochastic fit uncertainty but poor for recovering
instantaneous morphology.

\subsection{Noise-Blur Tradeoff}

Longer integration generally reduces stochastic measurement noise. A simple first-order rule is:
\begin{equation}
  \sigma_{\mathrm{stat}}(T_{\mathrm{int}})
  \propto
  \frac{1}{\sqrt{T_{\mathrm{int}}}}.
\end{equation}
But longer integration increases motion-induced model mismatch. A local first-order approximation
is:
\begin{equation}
  \sigma_{\mathrm{motion}}^2
  \approx
  \frac{T_{\mathrm{int}}^2}{12}
  \left(\mathbf{v}^{\top}\nabla u\right)^2.
\end{equation}
Therefore an effective uncertainty for fast-moving structures should include both:
\begin{equation}
  \sigma_{\mathrm{eff}}^2
  =
  \sigma_{\mathrm{fit}}^2
  +
  \sigma_{\mathrm{floor}}^2
  +
  \sigma_{\mathrm{motion}}^2
  +
  \sigma_{\mathrm{systematic}}^2.
\end{equation}
This expression should not be presented as final ISR theory yet. It is a practical synthetic and
diagnostic model that prevents long-integration products from being treated as overconfident point
measurements during active times.

\section{Recommended Synthetic Velocity Suite}

\subsection{Velocity Values}

The next synthetic suite should include:
\begin{equation}
  \|\mathbf{v}\| \in
  \{0,\;0.36,\;0.75,\;1.0,\;1.5,\;2.0,\;3.0\}~\kmps.
\end{equation}

Interpretation:
\begin{center}
\begin{tabular}{ll}
\toprule
Velocity & Purpose \\
\midrule
$0$ & static control \\
$0.36~\kmps$ & quiet/moderate baseline, current case \\
$0.75~\kmps$ & active but usually still trackable at short cadence \\
$1.0~\kmps$ & strong convection \\
$1.5~\kmps$ & active/storm-time stress \\
$2.0~\kmps$ & severe active-time case \\
$3.0~\kmps$ & out-of-distribution stress test \\
\bottomrule
\end{tabular}
\end{center}

These values are also physically interpretable from the $\mathbf{E}\times\mathbf{B}$ drift scale.
For a high-latitude magnetic field of order $B\sim 5\times 10^{-5}~\mathrm{T}$,
\begin{equation}
  v_E = \frac{E}{B}.
\end{equation}
Thus $E=50~\mathrm{mV\,m^{-1}}$ corresponds to approximately $1~\kmps$, and
$E=100~\mathrm{mV\,m^{-1}}$ corresponds to approximately $2~\kmps$. Such values are not
quiet-time values, but they are appropriate stress cases for active high-latitude convection
\cite{Kelley2009,Greenwald1995SuperDARN,Chisham2007SuperDARN}.

\subsection{Velocity Directions}

Do not use only one direction. Use a small set of directions:
\begin{equation}
  \phi \in \{0^\circ,\;45^\circ,\;90^\circ,\;135^\circ\},
\end{equation}
or align directions relative to the radar beam fan:
\begin{enumerate}[leftmargin=*]
  \item along the densest beam fan;
  \item across the beam fan;
  \item diagonal to the beam geometry;
  \item approximately field-aligned projection if magnetic coordinates are introduced.
\end{enumerate}

The same speed can be easy or hard depending on whether the structure moves along well-sampled
directions or through gaps between beams.

\section{Beam Count Strategy}

The current synthetic work includes 42-beam and 23-beam cases. An 11-beam case should be added,
but it should be interpreted correctly.

\begin{center}
\begin{tabular}{lll}
\toprule
Beam count & Role & Expected outcome \\
\midrule
42 & optimistic/method-development case & should reconstruct moderate motion well \\
23 & realistic sparse benchmark & main scientific benchmark \\
11 & stress/failure case & should expose low-reliability regions \\
\bottomrule
\end{tabular}
\end{center}

The 11-beam case may indeed be too sparse for strong reconstruction. That is not a reason to omit
it. The final tool must be able to tell a user when an interpolated field is not trustworthy. If an
11-beam high-speed case produces a visually smooth reconstruction but the reliability layer correctly
marks large regions as unreliable, that is a project success.

\section{Synthetic Experiment Matrix}

\subsection{Stage 1: Manageable Core Matrix}

A full factorial sweep over velocity, direction, integration time, beam count, architecture, and loss
weights would become too large. The first serious matrix should be:
\begin{enumerate}[leftmargin=*]
  \item beam counts: 42, 23, 11;
  \item velocities: $0.36$, $1.0$, $2.0~\kmps$;
  \item integration times: 1, 2, 5, 10 minutes;
  \item direction: one standard direction, preferably the current direction;
  \item regularization: data-only, moderate baseline, and selected stronger case.
\end{enumerate}

This is already:
\[
  3 \times 3 \times 4 \times 3 = 108
\]
training conditions if all are run. If runtime is high, reduce Stage 1 to:
\[
  \mathrm{beams}\in\{42,23,11\},\quad
  v\in\{0.36,1.0,2.0\},\quad
  T_{\mathrm{int}}\in\{1,5,10\},
\]
with only the moderate regularization setting and a data-only control for selected cases.

\subsection{Stage 2: Directional Stress}

After Stage 1, fix the most informative speeds, probably $1.0$ and $2.0~\kmps$, and test
several directions. This tells us whether performance is controlled by speed alone or by
speed relative to beam geometry.

\subsection{Stage 3: Velocity-Aware Models}

Compare:
\begin{enumerate}[leftmargin=*]
  \item standard INR:
  \[
    f_{\theta}(x,y,t);
  \]
  \item co-moving coordinate INR:
  \[
    f_{\theta}(x-u_x t,\; y-u_y t,\; t);
  \]
  \item weak advection diagnostic, not yet a training loss:
  \[
    r_{\mathrm{adv}} = f_t + u_x f_x + u_y f_y.
  \]
\end{enumerate}

The co-moving coordinate version is recommended before an advection residual loss because it is
still an interpolation representation rather than a PDE-constrained PINN. If an advection residual
is later used, it should be presented as a diagnostic or a soft transport prior, not as a full
physical model.

\section{Using Line-of-Sight Velocities}

Radar products include line-of-sight velocity. These should be used, but carefully.

\subsection{Synthetic Velocity Generation}

For synthetic data, define a true horizontal velocity:
\begin{equation}
  \mathbf{v}_{h} = (u_x,u_y).
\end{equation}
For beam unit vector $\hat{\mathbf{b}}_j$, the synthetic line-of-sight velocity is:
\begin{equation}
  \vlos^{(j)}
  =
  \mathbf{v}_{h}\cdot \hat{\mathbf{b}}_{h}^{(j)}
  +
  \epsilon_v,
\end{equation}
where $\hat{\mathbf{b}}_{h}^{(j)}$ is the horizontal projection of the beam direction. This allows
the synthetic benchmark to include realistic velocity observations without pretending that every
beam directly measures the full horizontal vector.

\subsection{First Real-Data Use}

For real data, the first use of $\vlos$ should be conservative:
\begin{enumerate}[leftmargin=*]
  \item use $\vlos$ as a speed/activity proxy;
  \item use it to flag intervals where long integration times may be motion blurred;
  \item use it as an input feature or reliability feature;
  \item use it to estimate $\sigma_{\mathrm{motion}}$;
  \item avoid claiming a full vector velocity field from monostatic LOS data without additional
  assumptions.
\end{enumerate}

Later, a smoothed horizontal velocity field can be estimated from multiple beams or external
context, but that should be a separate methodological step.

\section{ISR Uncertainty and Error Modeling}

\subsection{Formal Fit Errors}

The ISR error document emphasizes that reported variables such as $d\Ne$, $d\Te$, $d\Ti$, and
$d\Vi$ are formal fitting uncertainties, not total truth errors \cite{Diaz2026ISRError}. They are
derived from the nonlinear spectral fitting problem and reflect the local posterior covariance under
the assumed forward model \cite{Lehtinen1986,LehtinenHuuskonen1996,Virtanen2008,Virtanen2021}.

For a transformed target $u=\log_{10}\Ne$, the uncertainty should be transformed as:
\begin{equation}
  \sigma_u
  \approx
  \frac{\sigma_{\Ne}}{\ln(10)\Ne}.
\end{equation}

\subsection{Heteroscedastic Data Loss}

The data loss should become:
\begin{equation}
  \mathcal{L}_{\mathrm{data},\sigma}
  =
  \frac{1}{N}
  \sum_{i=1}^{N}
  \frac{
    \left(f_{\theta}(\mathbf{x}_i,t_i)-u_i\right)^2
  }{
    \sigma_{\mathrm{eff},i}^2
  }.
\end{equation}

The effective uncertainty should include a floor:
\begin{equation}
  \sigma_{\mathrm{eff},i}
  =
  \max(\sigma_{u,i},\sigma_{\mathrm{floor}}),
\end{equation}
or a richer expression:
\begin{equation}
  \sigma_{\mathrm{eff},i}^2
  =
  \sigma_{u,i}^2
  +
  \sigma_{\mathrm{floor}}^2
  +
  \sigma_{\mathrm{motion},i}^2
  +
  \sigma_{\mathrm{support},i}^2.
\end{equation}

\subsection{Normalized Residual Diagnostics}

Use:
\begin{equation}
  z_i =
  \frac{
    f_{\theta}(\mathbf{x}_i,t_i)-u_i
  }{
    \sigma_{\mathrm{eff},i}
  }.
\end{equation}
Then track:
\begin{equation}
  \chi^2_{\nu}
  \approx
  \frac{1}{N}
  \sum_i z_i^2.
\end{equation}
Values near one indicate rough consistency if the error model is appropriate. Values much larger
than one suggest under-estimated uncertainty or model mismatch. Values much smaller than one
suggest over-estimated uncertainty or over-smoothed predictions.

\subsection{Systematic Effects Not Captured by Formal Error Bars}

Formal fit uncertainties do not fully capture:
\begin{enumerate}[leftmargin=*]
  \item ion composition assumptions;
  \item non-Maxwellian distributions;
  \item beam-to-beam calibration bias;
  \item range-lag ambiguity and range mixing;
  \item velocity shear during integration;
  \item low-elevation gain/SNR degradation;
  \item I/Q and processing artifacts;
  \item temporal nonstationarity during integration.
\end{enumerate}
These should feed the reliability map even when they are not included directly in training.

\section{Anisotropy and the Vertical Dimension}

\subsection{Horizontal and Vertical Directions Should Differ}

The vertical coordinate $z$ should not be treated as simply another horizontal coordinate. The
ionosphere has strong vertical stratification, layer structure, scale-height behavior, composition
transitions, and range ambiguity. Horizontal interpolation across sparse beams and vertical
interpolation across range gates have different physical meanings.

For full 4D reconstruction, use:
\begin{equation}
  f_{\theta}(x,y,z,t)
  =
  b_{\phi}(z,t)
  +
  r_{\theta}(x,y,z,t),
\end{equation}
where $b_{\phi}$ is a smooth background vertical profile and $r_{\theta}$ is a residual 4D INR.
This decomposition gives the model a physically meaningful vertical baseline while allowing
horizontal structure.

\subsection{Anisotropic Curvature}

The 4D loss should generalize to:
\begin{equation}
  \mathcal{L}
  =
  \mathcal{L}_{\mathrm{data},\sigma}
  +
  \lambda_{xy}\mathcal{L}_{xy}
  +
  \lambda_z\mathcal{L}_z
  +
  \lambda_t\mathcal{L}_t.
\end{equation}
The vertical penalty can be:
\begin{equation}
  \mathcal{L}_z
  =
  \E[f_{zz}^2],
\end{equation}
but for sharp layers a robust penalty should be tested:
\begin{equation}
  \mathcal{L}_z^{\mathrm{Huber}}
  =
  \E\left[\rho_{\delta}(f_{zz})\right],
\end{equation}
where $\rho_{\delta}$ is the Huber loss. Cauchy or total-variation-like penalties are also worth
testing because full-profile ISR analysis and Bayesian inversion literature warns that stationary
Tikhonov-type smoothing can suppress real sharp structures \cite{Diaz2026ISRError,Hansen1992LCurve}.

\subsection{Magnetic Coordinates}

For high-latitude data, geographic $(x,y,z)$ may not always be the best coordinate system.
Potential future coordinate systems include:
\begin{enumerate}[leftmargin=*]
  \item geographic/cartesian coordinates;
  \item magnetic apex coordinates;
  \item coordinates decomposed into directions parallel and perpendicular to $\mathbf{B}$;
  \item beam/range/time coordinates for data-ingest diagnostics.
\end{enumerate}
The EISCAT 3D uncertainty literature emphasizes the importance of geometry, viewing angle relative
to the magnetic field, and experiment design \cite{Hatch2025E3D}.

\section{Architecture Roadmap}

\subsection{SIREN as the Main Baseline}

SIREN remains the correct baseline because:
\begin{enumerate}[leftmargin=*]
  \item it maps coordinates to continuous values;
  \item it can represent smooth fields and derivatives;
  \item it works naturally with derivative regularization;
  \item frequency scaling acts as a spectral prior.
\end{enumerate}
However, the SIREN frequency parameter $\omega_0$ must be treated as a hyperparameter and not
as a universal constant.

\subsection{FINER and FINER++}

FINER and FINER++ introduce variable-periodic activation functions to tune spectral support more
flexibly than standard SIREN \cite{Liu2023FINER,Zhu2024FINERPP}. These should be tested if:
\begin{enumerate}[leftmargin=*]
  \item low-frequency SIREN is too smooth;
  \item high-frequency SIREN creates ringing;
  \item one fixed $\omega_0$ cannot represent both smooth background and sharp structures.
\end{enumerate}

\subsection{K-Planes and Factorized 4D Representations}

K-Planes factorizes high-dimensional fields into learned planes and is explicitly designed to move
from static to dynamic scenes \cite{FridovichKeil2023KPlanes}. It is attractive for full
$(x,y,z,t)$ reconstruction because it can make dimension-specific priors easier to implement. For
example, temporal smoothness can be applied to time-containing planes and vertical priors can be
separated from horizontal ones.

\subsection{Hash or Grid Hybrid Representations}

Multiresolution hash encodings can drastically accelerate neural field training
\cite{Muller2022InstantNGP}. They should be considered if training cost becomes limiting.
However, they may make derivative priors less clean, so they should not replace SIREN before
the uncertainty-aware baseline is stable.

\subsection{Radar Neural Fields}

Radar-neural-field papers such as RF4D and SpINR are conceptually relevant because they show
that radar-like data can be represented by neural fields \cite{RF4D2024,SpINR2024}. However,
they are not directly equivalent to ISR plasma-parameter reconstruction. ISR products are fitted
plasma parameters with formal uncertainties, not raw FMCW point clouds or automotive radar
returns.

\section{Loss-Weight Selection Protocol}

\subsection{Why the Current Sweep Is Not Final}

The four-case sweep is scientifically useful but not sufficient as final evidence. It showed that
regularization matters and that stronger is not automatically better. A final study requires a
structured selection protocol.

\subsection{Recommended Weight Grid}

Use target ratios rather than raw lambdas:
\begin{equation}
  \rho_{xy}
  =
  \frac{\lambda_{xy}\mathcal{L}_{xy}}
  {\mathcal{L}_{\mathrm{data}}},
  \qquad
  \rho_t
  =
  \frac{\lambda_t\mathcal{L}_t}
  {\mathcal{L}_{\mathrm{data}}}.
\end{equation}

Recommended first grid:
\begin{equation}
  \rho_{xy},\rho_t \in
  \{0,\;0.05,\;0.1,\;0.3,\;0.7,\;1.0\}.
\end{equation}

Select weights by validation performance, not training loss. The primary selection metrics should
be withheld-beam error, withheld-time error, high-gradient error, and calibration quality.

\subsection{Adaptive Loss Balancing}

Compare the fixed-ratio method with:
\begin{enumerate}[leftmargin=*]
  \item L-curve analysis for the data-fit/regularization tradeoff
  \cite{Hansen1992LCurve};
  \item GradNorm-style gradient balancing \cite{Chen2018GradNorm};
  \item SoftAdapt-style loss-rate balancing \cite{Heydari2019SoftAdapt};
  \item uncertainty weighting for multi-task losses \cite{Kendall2018MultiTask};
  \item ReLoBRaLo-style balancing, used only as a multi-loss optimizer idea, not as a PINN
  claim \cite{Bischof2021ReLoBRaLo};
  \item gradient-flow diagnostics inspired by PINN loss pathology studies
  \cite{Wang2021GradientPathologies}.
\end{enumerate}

The most defensible phrasing is:
\begin{quote}
We seek the smoothest and most reliable field whose uncertainty-normalized data misfit remains
acceptable under withheld-beam and withheld-time validation.
\end{quote}

\section{Reliability Maps}

\subsection{Reliability Is a Deliverable, Not a Decoration}

The final tool should show both interpolation and reliability. A region marked unreliable is not a
failure; it is valuable scientific information. The reliability map should warn users when the model
is extrapolating across poorly observed space, long integration time, high motion, low SNR, or high
dropout/conformal uncertainty.

\subsection{MC Dropout Route}

MC dropout provides stochastic forward passes:
\begin{equation}
  \{\hat{u}^{(m)}(\mathbf{x},t)\}_{m=1}^{M}.
\end{equation}
From these compute:
\begin{align}
  \mu(\mathbf{x},t) &= \frac{1}{M}\sum_m \hat{u}^{(m)}(\mathbf{x},t),\\
  \sigma_{\mathrm{MC}}^2(\mathbf{x},t)
  &=
  \frac{1}{M-1}\sum_m
  \left(\hat{u}^{(m)}(\mathbf{x},t)-\mu(\mathbf{x},t)\right)^2.
\end{align}
This is an epistemic uncertainty proxy \cite{Gal2016Dropout,KendallGal2017}.

\subsection{Conformal Prediction Route}

Conformal prediction wraps a trained model with calibrated intervals. For calibration samples,
define a normalized score:
\begin{equation}
  s_i
  =
  \frac{
    |u_i-\hat{u}_i|
  }{
    a_i
  },
\end{equation}
where $a_i$ is a scale factor. For this project:
\begin{equation}
  a_i^2
  =
  \sigma_{\mathrm{ISR},i}^2
  +
  \sigma_{\mathrm{MC},i}^2
  +
  \sigma_{\mathrm{motion},i}^2
  +
  \sigma_{\mathrm{support},i}^2.
\end{equation}
Let $q_{1-\alpha}$ be the appropriate conformal quantile of the calibration scores. Then:
\begin{equation}
  \hat{u}(\mathbf{x},t)
  \pm
  q_{1-\alpha}a(\mathbf{x},t)
\end{equation}
is a calibrated prediction interval under the assumptions of the conformal split
\cite{AngelopoulosBates2021,Romano2019CQR}.

\subsection{Group-Conditioned Reliability}

A single global interval may be misleading. Use bins or groups:
\begin{enumerate}[leftmargin=*]
  \item altitude range;
  \item support distance;
  \item beam elevation;
  \item SNR or reported $\sigma$;
  \item integration time;
  \item velocity/motion-blur class;
  \item interior vs boundary regions.
\end{enumerate}
This is similar in spirit to Mondrian or group-conditioned conformal prediction, and it helps ensure
that high-confidence regions and weak-support regions are not averaged into one misleading
calibration score.

\subsection{Reliability Map Components}

The reliability map should combine:
\begin{enumerate}[leftmargin=*]
  \item ISR aleatoric uncertainty;
  \item MC dropout dispersion or conformal interval width;
  \item distance to nearest measurement;
  \item beam density/coverage;
  \item beam-angle diversity;
  \item integration-time motion blur;
  \item window-overlap disagreement;
  \item withheld-validation residual patterns;
  \item quality flags and systematic-risk indicators.
\end{enumerate}

\section{Validation Design}

\subsection{Synthetic Validation}

Synthetic validation should include:
\begin{enumerate}[leftmargin=*]
  \item dense truth error;
  \item gradient error;
  \item high-gradient region error;
  \item boundary vs interior error;
  \item near-observation vs far-support error;
  \item motion-speed and integration-time stratification;
  \item beam-count stratification.
\end{enumerate}

\subsection{Real-Data Validation}

Real-data validation should use:
\begin{enumerate}[leftmargin=*]
  \item random withheld points for debugging;
  \item withheld range gates;
  \item withheld altitude bands;
  \item withheld time blocks;
  \item withheld beams;
  \item combined withheld beam/time stress tests.
\end{enumerate}

\subsection{Metrics}

Report:
\begin{enumerate}[leftmargin=*]
  \item RMSE, MAE, bias, and 95th percentile absolute error;
  \item normalized residuals and reduced chi-square-like metrics;
  \item empirical coverage of nominal 50\%, 80\%, and 90\% intervals;
  \item interval width;
  \item coverage by altitude;
  \item coverage by support-distance bin;
  \item coverage by beam/elevation/SNR bin;
  \item coverage by integration time and velocity class.
\end{enumerate}

\section{Recommended Final Tool}

The final deliverable should be an interactive analysis tool, not only a training script. It should
show:
\begin{enumerate}[leftmargin=*]
  \item reconstructed $\log_{10}\Ne$ slices;
  \item later $\Te$ and $\Ti$ slices;
  \item observed beam/range samples overlaid;
  \item reliability map;
  \item interval-width map;
  \item support/observability map;
  \item integration-time or motion-blur warning layer;
  \item optional mask that hides low-reliability regions;
  \item calibration summary, e.g. ``nominal 90\% interval achieved 87\% coverage on withheld
  beams.''
\end{enumerate}

The most important user-interface principle is that unreliable regions should remain visible but
clearly marked. The scientific value is not only in showing the interpolation, but in showing where
the interpolation should not be trusted.

\section{Recommended Near-Term Tasks}

\subsection{Task 1: Add Integration-Aware Synthetic Generator}

Modify the synthetic observation generator so each observation can be either instantaneous or
integrated:
\begin{equation}
  y_i^{\mathrm{inst}} = u(\mathbf{r}_i,t_i),
  \qquad
  y_i^{\mathrm{int}} =
  \frac{1}{T_{\mathrm{int}}}\int u(\mathbf{r}_i,\tau)\dd\tau.
\end{equation}
Numerically approximate the integral by sub-sampling within the integration window.

\subsection{Task 2: Add Velocity and Integration Metadata}

Every synthetic output directory should record:
\begin{enumerate}[leftmargin=*]
  \item true velocity vector;
  \item speed;
  \item direction;
  \item integration time;
  \item cadence;
  \item beam count;
  \item feature width;
  \item $C_{\mathrm{samp}}$ and $C_{\mathrm{int}}$.
\end{enumerate}

\subsection{Task 3: Add 11-Beam Stress Geometry}

Use 11 beams as a reliability stress test. Do not expect uniform accuracy. Expect the reliability
map to degrade in unsampled regions.

\subsection{Task 4: Compare Instantaneous-Truth and Integrated-Truth Metrics}

For integrated products, evaluate against two truths:
\begin{enumerate}[leftmargin=*]
  \item instantaneous truth at the nominal timestamp;
  \item integrated truth over the same time window.
\end{enumerate}
If the model is trained on integrated observations, it may be unfair to demand exact instantaneous
truth recovery unless a deblurring or motion-compensation model is included.

\subsection{Task 5: Add Velocity-Aware Coordinates}

Test:
\begin{equation}
  (x,y,t)
  \quad \mathrm{vs.}\quad
  (x-u_x t,\;y-u_y t,\;t).
\end{equation}
This is the lowest-complexity way to see whether known or estimated velocity can improve
fast-moving reconstruction without turning the project into a PINN.

\subsection{Task 6: Add Reliability Prototype}

Start with:
\begin{equation}
  R(\mathbf{x},t)
  =
  g\left(
    d_{\mathrm{nearest}},
    \sigma_{\mathrm{ISR}},
    \sigma_{\mathrm{MC}},
    T_{\mathrm{int}},
    \|\mathbf{v}\|,
    \mathrm{beam\ density}
  \right),
\end{equation}
where $R$ is a reliability score and $g$ is initially a transparent hand-designed function. Later
replace or calibrate this with conformal intervals.

\section{Conclusion}

The project is no longer simply testing whether a SIREN can draw a plausible plasma-density
image. It is becoming a statistically and physically informed reconstruction framework for sparse
ISR products. The key scientific question is not merely whether the INR can interpolate. The key
question is:
\begin{quote}
Can the INR interpolate sparse ISR plasma fields while honestly indicating where the result is
data constrained and where it is not?
\end{quote}

The next stage should make velocity, integration time, beam sparsity, uncertainty, and reliability
central experimental axes. In particular, integration time is mandatory. At active-time velocities,
one-minute, two-minute, five-minute, and ten-minute products correspond to fundamentally
different inverse problems. A reliable reconstruction tool must know this, expose it, and warn the
user when long integration, sparse beams, or fast motion make the interpolation unreliable.

\begin{thebibliography}{99}

\bibitem{Diaz2026Proposal}
J. Diaz Pena.
\newblock \emph{Uncertainty-Aware AI to Reconstruct 4D Ionospheric Structure from Sparse Multi-Beam Radar Data}.
\newblock UAI PostDoc proposal, Google Drive project folder, 2026.
\newblock \url{https://drive.google.com/file/d/1k37yBuWw5WK526jeX9Cg0AD9l357yV3J/view}

\bibitem{Diaz2026ISRError}
J. Diaz Pena.
\newblock \emph{The Physics and Mathematics of Uncertainty in Incoherent Scatter Radar Data: An Analysis of AMISR Fitting Processes and Error Propagation}.
\newblock Internal project report, Google Drive project folder, 2026.
\newblock \url{https://drive.google.com/file/d/1USd56Wczouo_HAWEH3zIGyNd2jzOo24o/view}

\bibitem{Diaz2026INRReport}
J. Diaz Pena.
\newblock \emph{AMISR Implicit Neural Representation for Ionospheric Plasma Density Reconstruction}.
\newblock Internal project report, Google Drive project folder, 2026.
\newblock \url{https://drive.google.com/file/d/1MYmq-zf65s5ODCWezga1P4W2LERVEE0n/view}

\bibitem{Diaz2026INRProgress}
J. Diaz Pena.
\newblock \emph{Detailed Notes on the 3D INR Radar Case and the Synthetic-Data Validation Case}.
\newblock Internal project report, Google Drive project folder, 2026.
\newblock \url{https://drive.google.com/file/d/1bsz7WgimiaG3qVxcHwgKWQjK7Uu1BeaL/view}

\bibitem{Gordon1958}
W. E. Gordon.
\newblock Incoherent scattering of radio waves by free electrons with applications to space exploration by radar.
\newblock \emph{Proceedings of the IRE}, 46(11):1824--1829, 1958.

\bibitem{Bowles1958}
K. L. Bowles.
\newblock Observation of vertical-incidence scatter from the ionosphere at 41 Mc/sec.
\newblock \emph{Physical Review Letters}, 1(12):454--455, 1958.

\bibitem{Farley1969}
D. T. Farley.
\newblock Incoherent scatter correlation function measurements.
\newblock \emph{Radio Science}, 4(10):935--953, 1969.

\bibitem{Lehtinen1986}
M. S. Lehtinen.
\newblock \emph{Statistical Theory of Incoherent Scatter Radar Measurements}.
\newblock EISCAT Technical Note 86/45, 1986.

\bibitem{LehtinenHuuskonen1996}
M. S. Lehtinen and A. Huuskonen.
\newblock General incoherent scatter analysis and GUISDAP.
\newblock \emph{Journal of Atmospheric and Terrestrial Physics}, 58(1--4):435--452, 1996.

\bibitem{HuuskonenLehtinen1996}
A. Huuskonen and M. S. Lehtinen.
\newblock The accuracy of incoherent scatter measurements: error estimates valid for high signal levels.
\newblock \emph{Journal of Atmospheric and Terrestrial Physics}, 58(1--4):453--463, 1996.

\bibitem{KudekiMilla2011}
E. Kudeki and M. A. Milla.
\newblock Incoherent scatter spectral theories. Part I: A general framework and results for small magnetic aspect angles.
\newblock \emph{IEEE Transactions on Geoscience and Remote Sensing}, 49(1):315--328, 2011.

\bibitem{HeinselmanNicolls2008}
C. J. Heinselman and M. J. Nicolls.
\newblock A Bayesian approach to electric field and E-region neutral wind estimation with the Poker Flat Advanced Modular Incoherent Scatter Radar.
\newblock \emph{Radio Science}, 43, RS5013, 2008.

\bibitem{Kelly2009PFISR}
J. D. Kelly, C. J. Heinselman, J. F. Vickrey, and R. R. Vondrak.
\newblock Initial results from Poker Flat Incoherent Scatter Radar.
\newblock \emph{Journal of Atmospheric and Solar-Terrestrial Physics}, 71(6--7):635--646, 2009.

\bibitem{Virtanen2008}
I. I. Virtanen, M. S. Lehtinen, T. Nygren, M. Orispaa, and M. Vierinen.
\newblock A Bayesian approach to incoherent scatter radar analysis.
\newblock \emph{Annales Geophysicae}, 26:571--586, 2008.

\bibitem{Virtanen2021}
I. I. Virtanen, H. W. Tesfaw, L. Roininen, S. Lasanen, and A. Aikio.
\newblock Bayesian filtering in incoherent scatter plasma parameter fits.
\newblock \emph{Journal of Geophysical Research: Space Physics}, 126, e2020JA028700, 2021.

\bibitem{Hatch2025E3D}
S. M. Hatch, I. Virtanen, K. M. Laundal, H. W. Tesfaw, J. Vierinen, D. R. Huyghebaert, A. Spicher, and J. C. Hessen.
\newblock Toolkit for incoherent scatter radar experiment design and applications to EISCAT 3D.
\newblock \emph{Annales Geophysicae}, 43:633--649, 2025.
\newblock \url{https://doi.org/10.5194/angeo-43-633-2025}

\bibitem{Greenwald1995SuperDARN}
R. A. Greenwald et al.
\newblock DARN/SuperDARN: A global view of the dynamics of high-latitude convection.
\newblock \emph{Space Science Reviews}, 71:761--796, 1995.

\bibitem{Chisham2007SuperDARN}
G. Chisham et al.
\newblock A decade of the Super Dual Auroral Radar Network (SuperDARN): Scientific achievements, new techniques and future directions.
\newblock \emph{Surveys in Geophysics}, 28:33--109, 2007.

\bibitem{Kelley2009}
M. C. Kelley.
\newblock \emph{The Earth's Ionosphere: Plasma Physics and Electrodynamics}.
\newblock Academic Press, second edition, 2009.

\bibitem{Hansen1992LCurve}
P. C. Hansen.
\newblock Analysis of discrete ill-posed problems by means of the L-curve.
\newblock \emph{SIAM Review}, 34(4):561--580, 1992.

\bibitem{Sitzmann2020SIREN}
V. Sitzmann, J. N. P. Martel, A. W. Bergman, D. B. Lindell, and G. Wetzstein.
\newblock Implicit neural representations with periodic activation functions.
\newblock \emph{Advances in Neural Information Processing Systems}, 2020.
\newblock \url{https://arxiv.org/abs/2006.09661}

\bibitem{Tancik2020Fourier}
M. Tancik, P. P. Srinivasan, B. Mildenhall, S. Fridovich-Keil, N. Raghavan, U. Singhal, R. Ramamoorthi, J. Barron, and R. Ng.
\newblock Fourier features let networks learn high frequency functions in low dimensional domains.
\newblock \emph{Advances in Neural Information Processing Systems}, 2020.
\newblock \url{https://arxiv.org/abs/2006.10739}

\bibitem{Liu2023FINER}
Z. Liu, H. Zhu, Q. Zhang, J. Fu, W. Deng, Z. Ma, Y. Guo, and X. Cao.
\newblock FINER: Flexible spectral-bias tuning in implicit neural representation by variable-periodic activation functions.
\newblock arXiv:2312.02434, 2023.
\newblock \url{https://arxiv.org/abs/2312.02434}

\bibitem{Zhu2024FINERPP}
H. Zhu, Z. Liu, Q. Zhang, J. Fu, W. Deng, Z. Ma, Y. Guo, and X. Cao.
\newblock FINER++: Building a family of variable-periodic functions for activating implicit neural representation.
\newblock arXiv:2407.19434, 2024.
\newblock \url{https://arxiv.org/abs/2407.19434}

\bibitem{FridovichKeil2023KPlanes}
S. Fridovich-Keil, G. Meanti, F. Warburg, B. Recht, and A. Kanazawa.
\newblock K-Planes: Explicit radiance fields in space, time, and appearance.
\newblock \emph{CVPR}, 2023.
\newblock \url{https://arxiv.org/abs/2301.10241}

\bibitem{Muller2022InstantNGP}
T. Muller, A. Evans, C. Schied, and A. Keller.
\newblock Instant neural graphics primitives with a multiresolution hash encoding.
\newblock \emph{ACM Transactions on Graphics}, 41(4), 2022.
\newblock \url{https://arxiv.org/abs/2201.05989}

\bibitem{RF4D2024}
RF4D authors.
\newblock \emph{RF4D: Neural Radar Fields for Novel View Synthesis in Outdoor Dynamic Scenes}.
\newblock Project paper in UAI Drive paper folder.

\bibitem{SpINR2024}
SpINR authors.
\newblock \emph{SpINR: Neural Volumetric Reconstruction for FMCW Radars}.
\newblock Project paper in UAI Drive paper folder.

\bibitem{Gal2016Dropout}
Y. Gal and Z. Ghahramani.
\newblock Dropout as a Bayesian approximation: Representing model uncertainty in deep learning.
\newblock \emph{International Conference on Machine Learning}, 2016.
\newblock \url{https://arxiv.org/abs/1506.02142}

\bibitem{KendallGal2017}
A. Kendall and Y. Gal.
\newblock What uncertainties do we need in Bayesian deep learning for computer vision?
\newblock \emph{Advances in Neural Information Processing Systems}, 2017.
\newblock \url{https://arxiv.org/abs/1703.04977}

\bibitem{AngelopoulosBates2021}
A. N. Angelopoulos and S. Bates.
\newblock A gentle introduction to conformal prediction and distribution-free uncertainty quantification.
\newblock arXiv:2107.07511, 2021.
\newblock \url{https://arxiv.org/abs/2107.07511}

\bibitem{Romano2019CQR}
Y. Romano, E. Patterson, and E. J. Candes.
\newblock Conformalized quantile regression.
\newblock \emph{Advances in Neural Information Processing Systems}, 2019.
\newblock \url{https://arxiv.org/abs/1905.03222}

\bibitem{Chen2018GradNorm}
Z. Chen, V. Badrinarayanan, C.-Y. Lee, and A. Rabinovich.
\newblock GradNorm: Gradient normalization for adaptive loss balancing in deep multitask networks.
\newblock \emph{International Conference on Machine Learning}, 2018.
\newblock \url{https://arxiv.org/abs/1711.02257}

\bibitem{Heydari2019SoftAdapt}
A. A. Heydari, C. A. Thompson, and A. Mehmood.
\newblock SoftAdapt: Techniques for adaptive loss weighting of neural networks with multi-part loss functions.
\newblock arXiv:1912.12355, 2019.
\newblock \url{https://arxiv.org/abs/1912.12355}

\bibitem{Kendall2018MultiTask}
A. Kendall, Y. Gal, and R. Cipolla.
\newblock Multi-task learning using uncertainty to weigh losses for scene geometry and semantics.
\newblock \emph{CVPR}, 2018.
\newblock \url{https://arxiv.org/abs/1705.07115}

\bibitem{Wang2021GradientPathologies}
S. Wang, Y. Teng, and P. Perdikaris.
\newblock Understanding and mitigating gradient flow pathologies in physics-informed neural networks.
\newblock \emph{SIAM Journal on Scientific Computing}, 43(5):A3055--A3081, 2021.
\newblock \url{https://arxiv.org/abs/2001.04536}

\bibitem{Bischof2021ReLoBRaLo}
R. Bischof and M. A. Kraus.
\newblock Multi-objective loss balancing for physics-informed deep learning.
\newblock arXiv:2110.09813, 2021.
\newblock \url{https://arxiv.org/abs/2110.09813}

\end{thebibliography}

\end{document}
